\documentclass{beamer}

% Theme choice
\usetheme{Madrid} % You can change to e.g., Warsaw, Berlin, CambridgeUS, etc.

% Encoding and font
\usepackage[utf8]{inputenc}
\usepackage[T1]{fontenc}

% Graphics and tables
\usepackage{graphicx}
\usepackage{booktabs}

% Code listings
\usepackage{listings}
\lstset{
basicstyle=\ttfamily\small,
keywordstyle=\color{blue},
commentstyle=\color{gray},
stringstyle=\color{red},
breaklines=true,
frame=single
}

% Math packages
\usepackage{amsmath}
\usepackage{amssymb}

% Colors
\usepackage{xcolor}

% TikZ and PGFPlots
\usepackage{tikz}
\usepackage{pgfplots}
\pgfplotsset{compat=1.18}
\usetikzlibrary{positioning}

% Hyperlinks
\usepackage{hyperref}

% Title information
\title{Chapter 8: Midterm Review \& Exam}
\author{Teaching Assistant}
\institute{University Name}
\date{\today}

\begin{document}

\frame{\titlepage}

\begin{frame}[fragile]
    \frametitle{Chapter 8: Midterm Review \& Exam}
    
    \begin{block}{Welcome! Our Goal for Today}
        Welcome to our midterm review session! Today is all about connecting the dots.
        \vspace{0.5em}
        
        Our goal is to move beyond individual topics and \textbf{synthesize} the core programming concepts from the first seven weeks. By seeing how these ideas fit together, you'll be better prepared for the exam.
    \end{block}
    
    \vspace{1em}
    
    \textbf{Today's Agenda:}
    \begin{itemize}
        \item Exam Logistics \& Strategy
        \item Comprehensive Review of Weeks 1-7 Topics
        \item Integrated Practice Problem
        \item Open Q\&A Session
    \end{itemize}
    
    \vspace{\fill}
    
    \begin{alertblock}{Key Mindset for Today}
        Think \textbf{synthesis}, not just memorization. Focus on \textit{how} and \textit{why} concepts work together.
    \end{alertblock}
    
\end{frame}

\begin{frame}[fragile]
    \frametitle{Today's Agenda: Our Roadmap to Success}
    
    \begin{columns}[T] % Top-aligned columns
        \begin{column}{0.5\textwidth}
            \begin{block}{1. Exam Logistics \& Strategy}
                \begin{itemize}
                    \item \textbf{What:} We'll cover the format, timing, and tools (Google Colab).
                    \item \textbf{Why:} Knowing the "rules of the game" helps you focus on solving problems, not logistics.
                \end{itemize}
            \end{block}
            
            \vspace{1em}
            
            \begin{block}{2. Comprehensive Review}
                \begin{itemize}
                    \item \textbf{What:} A recap from fundamentals (variables) to more complex structures (loops, functions).
                    \item \textbf{Why:} Reinforce foundations and highlight the connections between topics.
                \end{itemize}
            \end{block}
        \end{column}
        
        \begin{column}{0.5\textwidth}
            \begin{block}{3. Integrated Practice Problem}
                \begin{itemize}
                    \item \textbf{What:} We will collaboratively solve a sample exam-style problem from start to finish.
                    \item \textbf{Why:} This is where theory meets practice in a realistic, multi-step context.
                \end{itemize}
            \end{block}
            
            \vspace{1em}
            
            \begin{block}{4. Open Q\&A Session}
                \begin{itemize}
                    \item \textbf{What:} Your dedicated time to ask anything. No question is too simple or complex.
                    \item \textbf{Why:} Your chance to clarify any lingering doubts before your final study push.
                \end{itemize}
            \end{block}
        \end{column}
    \end{columns}
    
\end{frame}

\begin{frame}[fragile]
    \frametitle{Exam Logistics \& Topics Covered}

    \begin{block}{Exam Format}
        \begin{itemize}
            \item \textbf{Where:} In-class, using Google Colab.
            \item \textbf{When:} During our regular class time (90 minutes).
            \item \textbf{Rules:} Open-note and open-course-materials.
            \item \textbf{\alert{Important:}} This is \textbf{not} an open-internet exam. Using communication tools, Stack Overflow, or generative AI is strictly prohibited.
        \end{itemize}
    \end{block}

    \begin{block}{Topics Covered (Weeks 1-7)}
        \begin{itemize}
            \item Variables, Data Types, and Operators
            \item Conditional Logic (if/elif/else)
            \item Lists (creation, indexing, methods)
            \item Iteration (for/while loops)
            \item Functions (definition, parameters, return)
            \item Problem-Solving \& Debugging
        \end{itemize}
    \end{block}
\end{frame}

\begin{frame}[fragile]
    \frametitle{Exam Strategy: How to Succeed}

    \begin{block}{A 4-Step Strategy for Success}
        \begin{enumerate}
            \item \textbf{Read All Questions First (2-3 mins)} \\
            Quickly skim the entire exam to gauge difficulty and plan your approach. This helps you budget your time effectively.

            \item \textbf{Solve Easiest Problems First} \\
            Start with your "quick wins" to build confidence and secure points early. Don't leave easy points on the table.

            \item \textbf{Don't Get Stuck (The 15-Minute Rule)} \\
            If you make no progress on a problem after 10-15 minutes, leave a comment and move on. You can always come back later.

            \item \textbf{Save Time to Review (5-10 mins)} \\
            Finish early to review your work. Run all cells from the top, check for simple errors, and ensure you met all requirements.
        \end{enumerate}
    \end{block}
\end{frame}

\begin{frame}[fragile]
    \frametitle{Core Concepts: Data \& Expressions - Part 1}
    \framesubtitle{Variables \& Data Types}
    
    A review of the fundamental building blocks of any program. Variables act as named containers for storing data.
    
    \begin{block}{Common Data Types}
        Different types of data require different variable types.
        \begin{itemize}
            \item \texttt{name = "Python"} \hfill \textit{(string: for text)}
            \item \texttt{version = 3} \hfill \textit{(integer: for whole numbers)}
            \item \texttt{pi\_approx = 3.14} \hfill \textit{(float: for decimal numbers)}
            \item \texttt{is\_beginner = True} \hfill \textit{(boolean: for True/False)}
        \end{itemize}
    \end{block}
    
    Each variable holds a value of a specific type, which determines the operations you can perform on it.
\end{frame}

\begin{frame}[fragile]
    \frametitle{Core Concepts: Data \& Expressions - Part 2}
    \framesubtitle{Operators \& User Input}
    
    Programs interact with users and process data using operators and functions.
    
    \begin{block}{Example: Calculating Age Next Year}
    The \texttt{input()} function reads user data as a string. We must convert this string to a number for calculations—a process called \textbf{type casting}.
    \begin{lstlisting}[language=Python]
% Get user input (which is always a string)
age_str = input("Enter your age: ")

% Type cast the string to an integer for math
age_int = int(age_str)

% Perform calculation and print the result
% Note: We must cast the number back to a string to join it
print("Next year you will be: " + str(age_int + 1))
    \end{lstlisting}
    \end{block}
    
    \textbf{Key Steps:}
    \begin{enumerate}
        \item Get data from the user with \texttt{input()}.
        \item Convert data type for calculation (\texttt{string} $\rightarrow$ \texttt{integer}).
        \item Use an operator (\texttt{+}) to compute the new value.
        \item Convert the result back to a string for printing with other text.
    \end{enumerate}
\end{frame}

\begin{frame}[fragile]
    \frametitle{Control Flow Part 1: Conditionals}
    Making decisions in your code with \texttt{if}, \texttt{elif}, and \texttt{else}. The program's path is determined by Boolean expressions.
    \vspace{1em}

    \begin{columns}[T]
        \begin{column}{0.45\textwidth}
            \begin{itemize}
                \item \textbf{\texttt{if}}: The entry point. Checks a condition. If \texttt{True}, its code block runs.
                \vspace{0.5em}
                \item \textbf{\texttt{elif}}: Optional. Checks a new condition \textit{only if} the previous \texttt{if}/\texttt{elif} was \texttt{False}.
                \vspace{0.5em}
                \item \textbf{\texttt{else}}: Optional. A "catch-all" that runs \textit{only if} all preceding conditions were \texttt{False}.
            \end{itemize}
        \end{column}
        
        \begin{column}{0.5\textwidth}
            \textbf{Example:}
            \begin{verbatim}
score = 75

if score >= 90:
    print("Excellent!")
elif score >= 60:
    print("Satisfactory.")
else:
    print("Needs improvement.")
            \end{verbatim}
        \end{column}
    \end{columns}

    \begin{block}{Key Idea}
        For any single \texttt{if-elif-else} chain, Python executes conditions top-to-bottom and runs \textbf{at most one} code block.
    \end{block}
\end{frame}

\begin{frame}[fragile]
    \frametitle{Conditionals: Execution Walkthrough}
    Let's trace the execution for \texttt{score = 75}. As soon as a condition is met, the rest of the chain is skipped.

    \begin{block}{Example Trace: \texttt{score = 75}}
        \begin{enumerate}
            \item \textbf{Check the \texttt{if} condition:} \\
            Is \texttt{75 >= 90}? \hfill \textbf{False}. \\
            \textit{Action: Skip this block and move on.}
            \begin{verbatim}
if score >= 90:
    print("Excellent!")
            \end{verbatim}
            
            \item \textbf{Check the \texttt{elif} condition:} \\
            Is \texttt{75 >= 60}? \hfill \textbf{True}. \\
            \textit{Action: Execute this block. Prints "Satisfactory."}
            \begin{verbatim}
elif score >= 60:
    print("Satisfactory.")
            \end{verbatim}
            
            \item \textbf{Exit the chain:} \\
            Because the \texttt{elif} was true, the \texttt{else} block is completely skipped.
        \end{enumerate}
    \end{block}
\end{frame}

\begin{frame}[fragile]
    \frametitle{Conditionals: A Common Pitfall}
    An \texttt{if-elif-else} chain is very different from a series of independent \texttt{if} statements.
    
    \begin{columns}[T]
        \begin{column}{0.5\textwidth}
            \begin{block}{\texttt{if-elif} Chain (Mutually Exclusive)}
                Only \textbf{one} block can possibly run.
\begin{verbatim}
score = 95
if score >= 90:
    print("Excellent!")
elif score >= 60:
    print("Satisfactory.")
\end{verbatim}
                \textbf{Output:} \texttt{Excellent!}
            \end{block}
        \end{column}
        \begin{column}{0.5\textwidth}
            \begin{block}{Multiple \texttt{if}s (Independent Checks)}
                \textbf{Each} \texttt{if} is tested separately. Multiple blocks could run.
\begin{verbatim}
score = 95
if score >= 90:
    print("Excellent!")
if score >= 60:
    print("Satisfactory.")
\end{verbatim}
                \textbf{Output:} Both messages print.
            \end{block}
        \end{column}
    \end{columns}
    
    \vspace{1em}
    \textbf{Key Takeaway:} Use an \texttt{if-elif} chain when you want to choose exactly one option from a set of possibilities.
\end{frame}

\begin{frame}[fragile]
    \frametitle{Control Flow Part 2: Iteration}
    \framesubtitle{For Loops (Definite Iteration)}

    Loops automate repetitive tasks, following the \textbf{DRY (Don't Repeat Yourself)} principle. Instead of writing code multiple times, we loop over it.

    \begin{block}{For Loops: Iterating over a sequence}
        Use a \texttt{for} loop when you know how many times you need to iterate, e.g., for every item in a collection. It's like going through a checklist.
        \vspace{0.5em}
        
        \textbf{Example 1: Looping over a list}
        \begin{verbatim}
# The variable `num` takes on each value from the list
numbers = [10, 20, 30]
for num in numbers:
    print(f"Current number is: {num}")
        \end{verbatim}

        \textbf{Example 2: Looping a fixed number of times}
        \begin{verbatim}
# range(3) generates a sequence: 0, 1, 2
for i in range(3):
    print("Hello!") # This line runs 3 times
        \end{verbatim}
    \end{block}
\end{frame}

\begin{frame}[fragile]
    \frametitle{Control Flow Part 2: Iteration}
    \framesubtitle{While Loops (Indefinite Iteration)}

    \begin{block}{While Loops: Repeating until a condition is false}
        Use a \texttt{while} loop when the number of repetitions is unknown and depends on a condition. The loop continues as long as the condition is \texttt{True}.
        
        \begin{enumerate}
            \item \textbf{Initialization:} Set up a variable before the loop.
            \item \textbf{Condition:} The expression checked before each iteration.
            \item \textbf{Update:} A change inside the loop to eventually make the condition false.
        \end{enumerate}
        
        \begin{verbatim}
count = 5           # 1. Initialization
while count > 0:    # 2. Condition
    print(count)
    count -= 1      # 3. Update Step
print("Blastoff!")
        \end{verbatim}
    \end{block}

    \begin{alertblock}{Critical Warning: Avoid Infinite Loops!}
        If you forget the update step (\texttt{count -= 1}), the condition will always be true, and the loop will run forever! Always ensure something inside the loop works toward stopping it.
    \end{alertblock}
\end{frame}

\begin{frame}[fragile]
    \frametitle{Data Structures: Lists (Part 1)}
    \framesubtitle{Definition \& Creation}

    A list is an \textbf{ordered} (maintains sequence) and \textbf{mutable} (changeable) collection of items, defined with square brackets \texttt{[]}.

    \begin{block}{Creating and Inspecting a List}
        \begin{verbatim}
# A list of strings
fruits = ["apple", "banana", "cherry"]

# A list can hold different data types
mixed_data = [10, "hello", 3.14, True]

# Get the number of items in the list
num_fruits = len(fruits) # Result is 3
        \end{verbatim}
    \end{block}
\end{frame}

\begin{frame}[fragile]
    \frametitle{Data Structures: Lists (Part 2)}
    \framesubtitle{Accessing \& Modifying Elements}

    We use an item's \textbf{index} (its position) to work with it.
    \begin{alertblock}{Key Point}
        Python uses \textbf{zero-based indexing}, meaning the first item is at index \texttt{0}.
    \end{alertblock}

    \begin{block}{Index-Based Operations}
        \begin{verbatim}
#          index:  0        1         2
fruits = ["apple", "banana", "cherry"]

# --- Accessing ---
# Get the first item (index 0)
first_fruit = fruits[0]      # "apple"

# Get the last item (negative indexing)
last_fruit = fruits[-1]     # "cherry"

# --- Modifying ---
# Change the item at index 1
fruits[1] = "blueberry" 
# fruits is now: ["apple", "blueberry", "cherry"]
        \end{verbatim}
    \end{block}
\end{frame}

\begin{frame}[fragile]
    \frametitle{Data Structures: Lists (Part 3)}
    \framesubtitle{Common Methods \& Slicing}

    \begin{block}{Common Methods: Changing List Size}
        Methods are functions that belong to the list object, called with dot notation.
        \begin{itemize}
            \item \texttt{.append(item)}: Adds a single item to the \textbf{end} of the list.
            \item \texttt{.pop(index)}: Removes the item at the specified index and \textbf{returns} it. If no index is given, it removes the last item.
        \end{itemize}
        \begin{verbatim}
# Original: ["apple", "blueberry", "cherry"]
fruits.append("orange")
# Now: ["apple", "blueberry", "cherry", "orange"]

removed_fruit = fruits.pop(1) # Removes "blueberry"
# removed_fruit is now: "blueberry"
        \end{verbatim}
    \end{block}

    \begin{block}{Slicing: Getting a Sub-List}
        You can extract a portion of a list using the slice syntax \texttt{[start:stop]}. The slice includes the \texttt{start} index but \textbf{excludes} the \texttt{stop} index.
        \begin{verbatim}
#          index:    0         1        2        3
all_items = ["pen", "paper", "stapler", "tape"]

# Get items from index 1 up to (but not including) 3
middle_items = all_items[1:3]  # ["paper", "stapler"]

# Get the first two items (from start up to index 2)
first_two = all_items[:2]      # ["pen", "paper"]
        \end{verbatim}
    \end{block}
\end{frame}

\begin{frame}[fragile]
    \frametitle{Modular Code: Functions}
    
    Functions are reusable, named blocks of code that perform a specific task. They are the primary way to structure code into logical, manageable units.
    
    \begin{block}{The DRY Principle: Don't Repeat Yourself}
        By defining code in one place and calling it when needed, you avoid duplication. If the logic needs to change, you only have to update it in one location.
    \end{block}
    
    \textbf{Why are functions essential?}
    \begin{itemize}
        \item \textbf{Organization:} Break down complex problems into smaller, logical pieces.
        \item \textbf{Reusability:} Write code once and use it many times.
        \item \textbf{Readability:} A well-named function like \texttt{calculate\_gross\_pay()} is much clearer than a block of raw calculations.
    \end{itemize}
    
\end{frame}

\begin{frame}[fragile]
    \frametitle{Anatomy of a Function}
    
    Let's break down the structure of a function definition and call.
    
    \begin{columns}[T]
        \begin{column}{0.6\textwidth}
            \begin{lstlisting}[language=Python]
# A function to calculate the area of a circle
def calculate_area(radius):  # Definition
    pi = 3.14159
    area = pi * (radius ** 2)
    return area              # Output

# Calling the function
circle_area = calculate_area(5) # 5 is the argument
            \end{lstlisting}
        \end{column}
        \begin{column}{0.4\textwidth}
            \begin{block}{Key Parts}
                \begin{itemize}
                    \item \texttt{def}: Keyword to \textbf{define} the function.
                    \item \texttt{radius}: A \textbf{parameter} (the expected input).
                    \item \texttt{return}: Keyword to send a value back (the \textbf{output}).
                    \item \texttt{5}: An \textbf{argument} (the actual value passed to the parameter).
                \end{itemize}
            \end{block}
        \end{column}
    \end{columns}

    \vspace{1em}
    \textbf{Mental Model:} Think of a function as a machine: \\
    \textbf{Input} (Parameters) $\rightarrow$ \textbf{Process} (Body) $\rightarrow$ \textbf{Output} (Return Value)
    
\end{frame}

\begin{frame}[fragile]
    \frametitle{The Art of Problem-Solving \& Debugging}
    The most critical skill is not just writing code, but solving problems systematically. The best programmers are systematic thinkers.
    \vspace{1em}

    \begin{block}{The Four-Step Process}
        This framework turns any complex problem into a manageable, solvable task.
        \begin{enumerate}
            \item \textbf{Understand the Problem:} Clearly identify inputs, outputs, and constraints.
            \item \textbf{Plan Your Algorithm:} Write down logical steps in pseudocode before touching the keyboard.
            \item \textbf{Translate to Code:} Implement your planned steps in Python.
            \item \textbf{Test \& Debug:} Run code with test cases. Use the traceback message to locate and fix errors.
        \end{enumerate}
    \end{block}

    \vspace{1em}
    \textit{We will explore these steps with an example: "Find the largest number in a list."}
\end{frame}

\begin{frame}[fragile]
    \frametitle{Steps 1 \& 2: Understand and Plan}
    \begin{block}{Example: Find the largest number in a list}
        \textbf{1. Understand the Problem}
        \begin{itemize}
            \item \textbf{Inputs:} A list of numbers (e.g., \texttt{[10, 5, 25, 8]}).
            \item \textbf{Outputs:} A single number, the maximum value (e.g., \texttt{25}).
            \item \textbf{Constraints/Edge Cases:} What if the list is empty? A safe choice is to return \texttt{None}.
        \end{itemize}
        \vspace{1em}
        \textbf{2. Plan Your Algorithm (Pseudocode)}
        \begin{lstlisting}[basicstyle=\ttfamily\scriptsize, frame=single]
FUNCTION find_max(numbers_list):
  // Handle the edge case first
  IF the list is empty:
    RETURN None

  // Assume the first number is the largest to start
  SET max_so_far = the first item in the list

  // Loop through the rest of the list
  FOR each number from the second item onwards:
    IF the current number > max_so_far:
      UPDATE max_so_far = the current number

  // Return the result
  RETURN max_so_far
        \end{lstlisting}
    \end{block}
\end{frame}

\begin{frame}[fragile]
    \frametitle{Steps 3 \& 4: Code and Debug}
    \textbf{3. Translate to Code}
    \begin{lstlisting}[language=Python]
def find_max(numbers_list):
    # Handle the edge case from our plan
    if not numbers_list:
        return None

    # Initialize max_so_far
    max_so_far = numbers_list[0]

    # Loop and compare (slice to start from the second item)
    for number in numbers_list[1:]:
        if number > max_so_far:
            max_so_far = number

    return max_so_far
    \end{lstlisting}

    \textbf{4. Test \& Debug} \\
    A traceback message is your best friend. If we forget the empty list check, we get:
    \begin{lstlisting}[language=bash, basicstyle=\ttfamily\tiny]
Traceback (most recent call last):
  File "main.py", line 6, in find_max
    max_so_far = numbers_list[0]
IndexError: list index out of range
    \end{lstlisting}
    This tells us exactly what (\texttt{IndexError}) and where (line 6) the problem is.

    \begin{alertblock}{Key Takeaway}
        \textbf{Slow down to speed up.} Time spent planning is paid back tenfold during coding and debugging.
    \end{alertblock}
\end{frame}

\begin{frame}[fragile]
    \frametitle{Integrated Practice Problem: Problem \& Plan (1/3)}
    Let's apply our problem-solving process to a sample exam-style question.

    \begin{block}{1. Understand the Problem}
        \begin{itemize}
            \item \textbf{Inputs:} A list of strings (e.g., \texttt{['apple', 'banana', 'kiwi']}) and an integer \texttt{min\_length} (e.g., \texttt{5}).
            \item \textbf{Output:} A \textbf{new list} containing only the words from the original that have a length greater than or equal to \texttt{min\_length}.
            \item \textbf{Constraint:} Do not modify the original list.
        \end{itemize}
    \end{block}

    \begin{block}{2. Plan (Pseudocode)}
        \begin{itemize}
            \item Define a function \texttt{filter\_long\_words} that accepts a list and a number.
            \item Create a new, empty list called \texttt{result\_list}.
            \item Loop through each \texttt{word} in the input list.
            \item Inside the loop, check if \texttt{len(word) >= min\_length}.
            \item If true, append the \texttt{word} to \texttt{result\_list}.
            \item After the loop finishes, return \texttt{result\_list}.
        \end{itemize}
    \end{block}
\end{frame}

\begin{frame}[fragile]
    \frametitle{Integrated Practice Problem: Implementation (2/3)}

    \begin{block}{3. Translate to Code (Python Implementation)}
        Here's how our pseudocode translates directly into a Python function. Each comment links back to a step in our plan.
        \begin{lstlisting}[language=Python, basicstyle=\small\ttfamily, commentstyle=\color{gray}]
def filter_long_words(word_list, min_length):
    # Create a new, empty list
    result_list = []
    
    # Loop through each word in the input list
    for word in word_list:
        # Inside the loop, check the condition
        if len(word) >= min_length:
            # If it is, append the word to our new list
            result_list.append(word)
            
    # After the loop, return the new list
    return result_list
        \end{lstlisting}
    \end{block}
\end{frame}

\begin{frame}[fragile]
    \frametitle{Integrated Practice Problem: Verification \& Review (3/3)}

    \begin{block}{4. Test \& Verify}
        We must test our function to ensure it works as expected.
        \begin{itemize}
            \item \textbf{Test Case 1: Standard Operation}
            \begin{lstlisting}[language=Python, basicstyle=\small\ttfamily]
words = ["python", "is", "fun", "and", "powerful"]
long_words = filter_long_words(words, 5)
print(long_words)
# Expected Output: ['python', 'powerful']
            \end{lstlisting}
            \item \textbf{Test Case 2: Edge Case (No words are long enough)}
            \begin{lstlisting}[language=Python, basicstyle=\small\ttfamily]
long_words_edge = filter_long_words(words, 10)
print(long_words_edge)
# Expected Output: []
            \end{lstlisting}
        \end{itemize}
    \end{block}

    \begin{block}{Key Concepts Reviewed}
        \begin{itemize}
            \item \textbf{Function Definition:} \texttt{def}, parameters, and \texttt{return}.
            \item \textbf{List Iteration:} Using a \texttt{for} loop.
            \item \textbf{Conditional Logic:} Using an \texttt{if} statement.
            \item \textbf{Built-in Functions:} Using \texttt{len()}.
            \item \textbf{List Manipulation:} Creating lists (\texttt{[]}) and using \texttt{.append()}.
        \end{itemize}
    \end{block}
\end{frame}

\begin{frame}[fragile]
    \frametitle{Q\&A and Final Reminders - Open Discussion}
    \begin{center}
        \Large \textbf{What concepts are you still unsure about?}
    \end{center}
    \vspace{1em}
    Let's discuss common sticking points. Do any of these spark a question?
    \begin{itemize}
        \item \textbf{Functions:} What is the difference between \texttt{print} and \texttt{return}?
        \item \textbf{Mutability:} Why can we change lists but not strings?
        \item \textbf{Looping Logic:} When to use a \texttt{for} loop vs. a \texttt{while} loop?
        \item \textbf{Data Structures:} When is a dictionary better than a list?
    \end{itemize}

    \begin{block}{Key Difference: \texttt{print} vs. \texttt{return}}
        \texttt{print()} displays a value for humans. The value cannot be used later. \\
        \texttt{return} sends a value back from a function to be stored and reused.
        \begin{lstlisting}[language=Python]
% This function's result CAN be stored
def return_sum(a, b):
    return a + b

result = return_sum(3, 5) # result is now 8
print(result * 2)          # Prints 16. We can use the returned value!
        \end{lstlisting}
    \end{block}
\end{frame}

\begin{frame}[fragile]
    \frametitle{Q\&A and Final Reminders - How to Prepare}
    \textbf{1. Redo Labs \& Assignments (Active Recall)}
    \begin{itemize}
        \item This is the most effective preparation. The exam is about problem-solving.
        \item \textbf{Method:} Start with a blank file. Try to solve the problem from scratch without looking at your old solution. This forces you to recall the entire process.
    \end{itemize}
    \vspace{1em}
    \textbf{2. Review Notes \& Readings (Synthesize Knowledge)}
    \begin{itemize}
        \item Connect ideas: How do loops, lists, and file I/O work together?
        \item \textbf{Pro-Tip:} Create a one-page summary sheet of key syntax. The act of making it is a powerful study tool.
    \end{itemize}
    \vspace{1em}
    \textbf{3. Use Course Resources}
    \begin{itemize}
        \item Don't study in isolation. Use the Q\&A forum (Ed/Piazza).
        \item Ask specific questions with code examples for the best help.
        \item Attend instructor and TA office hours.
    \end{itemize}
\end{frame}

\begin{frame}[fragile]
    \frametitle{Q\&A and Final Reminders - Final Advice}
    Your mental state is as important as your knowledge.

    \begin{block}{Before the Exam}
        \textbf{Get a good night's sleep!} A well-rested brain is far more effective at problem-solving than a tired one. All-night cramming is counterproductive.
    \end{block}
    
    \textbf{During the Exam:}
    \begin{itemize}
        \item Read every instruction carefully.
        \item If you get stuck on a problem, move on and come back later.
        \item It's better to get partial credit on all questions than a perfect score on one and leave others blank.
        \item Write pseudocode or comments to explain your logic if you can't finish the code; this can often earn partial credit.
    \end{itemize}
    \vspace{1em}
    You have put in the work and are well-prepared for this. Trust your preparation!
    
    \begin{center}
        \Large\textbf{Good luck!}
    \end{center}
\end{frame}


\end{document}